\documentclass[11pt,a4paper]{article}

\usepackage[utf8]{inputenc}
\usepackage[T1]{fontenc}
\usepackage[english]{babel}
\usepackage{amsmath, amssymb}
\usepackage{geometry}
\usepackage{booktabs}
\usepackage{array}
\usepackage{graphicx}
\usepackage{hyperref}
\usepackage{xcolor}
\usepackage{enumitem}
\usepackage{caption}

\geometry{margin=2.5cm}
\hypersetup{colorlinks=true, linkcolor=blue, urlcolor=blue}

\title{Presentation companion: the full narrative of the IFC\\
  predictive modeling project}
\author{Applied Statistics Project --- Presentation Preparation Notes}
\date{}

\begin{document}
\maketitle

\begin{abstract}
\noindent
This document is not a technical report: it is a speaking aid. It
walks through the entire project in the order it should be told out
loud, explains every number that appears on the poster, states
explicitly what was deliberately left out and why, and closes with
an anticipated question-and-answer section so that no reasonable
follow-up question catches the presenter unprepared. Read sections
1--7 to rehearse the talk; read section 8 the night before to
rehearse the defense.
\end{abstract}

\section{The one-sentence thesis}

\textbf{Italian municipal fragility is predictable and readable from
a handful of socio-economic indicators, but a residual geographic
component shows that the phenomenon is not purely individual --- it
is also spatial.}

Everything else in the project supports this sentence. If only one
sentence survives from the whole presentation, it should be this
one.

\section{The research question and why it matters}

The Italian Fragility Composite Index (IFC) is an official
municipal-level index. The project asks two questions:

\begin{enumerate}[leftmargin=2em]
  \item Can the IFC be statistically explained / predicted from
  independent socio-economic indicators (the GRINS V3 panel)?
  \item If so, what actually drives municipal fragility, and is the
  explanation purely socio-economic or does geography play an
  independent role?
\end{enumerate}

This matters because if a small set of readable indicators can
reconstruct most of the variance in an official fragility index,
that index can be monitored cheaply between official releases, and
the drivers identified become candidate levers for policy.

\section{The data, in one paragraph}

\textbf{Target}: raw IFC values for 2019 and 2021, from
\texttt{IFC\_Final\_Analysis\_Sorted.xlsx}; matches the official MFI
deciles 95.6\% of the time, confirming it is a faithful proxy for
the index actually used by policymakers.
\textbf{Predictors}: the GRINS V3 municipal panel, 549 raw columns,
engineered into 326 candidate covariates (counts normalised by
population plus a log1p transform, negative counts normalised by
population only, size variables log-transformed, rates left as is,
everything winsorised at the 1st/99th percentile to limit the
influence of outlier municipalities).
\textbf{Sample}: 15{,}782 observations, 7{,}886 unique
municipalities, two years (2019, 2021). Cross-validation is always
done by splitting whole municipalities into folds, never by row ---
splitting by row would leak the 2019 observation of a municipality
into the training fold while its 2021 twin sits in the test fold,
silently inflating R$^2$.

\section{The modeling pipeline, told as a sequence of questions}

Each step exists because it answers a question the previous step
left open. Tell it in this order.

\subsection*{Step 1 --- How much can be explained at all?}
A LASSO regression (10-fold CV, $\lambda_{\text{1se}}$) is run on
all 326 candidate covariates. It selects 191 of them. Out-of-sample
R$^2 = 0.825$ (repeated 5$\times$5 panel-aware CV, 25 folds). This is
the \emph{benchmark}: the ceiling of what a linear model can explain
with this feature set. The instructor's first reaction was correct:
191 covariates cannot be discussed one by one in any presentation.

\subsection*{Step 2 --- Can the same result be told with fewer
variables?}
An iterative pipeline reduces 191 covariates to a parsimonious set
of 25, explicit and reproducible:
\begin{enumerate}[leftmargin=2em]
  \item \textbf{VIF pruning}: drop the covariate with the highest
  Variance Inflation Factor, recompute, repeat, until all remaining
  VIFs are below 10 (34 covariates dropped this way).
  \item \textbf{Top-N ranking}: order the survivors by standardised
  effect size $|\beta \cdot \mathrm{SD}|$ and keep the top 25 ---
  this is what makes the variables comparable to each other despite
  living on different natural scales.
  \item \textbf{Significance filter}: any covariate that turns out
  non-significant once the other 24 are in the model is
  automatically replaced by the next candidate in the ranking,
  repeated until all 25 are significant.
\end{enumerate}
One explicit substantive decision precedes all three steps: the raw
municipal surface area is excluded upfront, on direct instructor
request, because once every other covariate is expressed as a
per-capita density or rate, a static geographic descriptor like
total surface area is not itself a fragility driver --- it is a
scale variable that the rest of the pipeline already controls for
implicitly.

Result: 25 covariates, every one significant at $p < 10^{-6}$,
maximum VIF $= 6.4$. Out-of-sample R$^2 = 0.784$ --- a controlled
loss of only $0.041$ against the 191-covariate benchmark, for an
87\% reduction in model size.

\subsection*{Step 3 --- Does geography add anything the 25 covariates
miss?}
A linear mixed model adds a random intercept structure on top of the
same 25 fixed effects:
$$
\text{IFC}_i = \beta_0 + \sum_{j=1}^{25}\beta_j X_{ji} +
u_{r(i)} + v_{p(i)\mid r(i)} + \varepsilon_i,
$$
$$
u_r \sim \mathcal{N}(0,\sigma_R^2), \qquad
v_{p\mid r} \sim \mathcal{N}(0,\sigma_P^2),
$$
with $r$ indexing 20 regions and $p$ indexing 107 provinces nested
inside regions. Four variants were fit and compared (region only,
region/province nested, taxonomy-based grouping, and the combination
of geography and taxonomy); the nested geographic model is the best
of the four. Out-of-sample R$^2 = 0.818$ --- within $0.007$ of the
191-covariate benchmark, recovering almost all of the accuracy lost
by going parsimonious. The intra-class correlation attributable to
geography is approximately 36\%: more than a third of the variance
left over after the 25 fixed effects is explained by which region
and province a municipality sits in, not by any covariate value.

\subsection*{Step 4 --- Is that geographic signal already exhausted,
or is something still missing?}
Global Moran's $I$ on the nested model's residuals, using a
symmetrised 5-nearest-neighbour spatial weight matrix:
$$ I = 0.234, \quad p < 0.001. $$
This is positive and highly significant: \emph{after} controlling
for 25 socio-economic covariates and region/province random
intercepts, neighbouring municipalities still share unexplained
fragility. Geography is not fully absorbed by administrative
boundaries. A teammate's local LISA decomposition on the same
residuals, same weights, independently recovers $I = 0.236$ ---
confirming the global number --- and locates it: 598 municipalities
form significant "High-High" clusters (under-predicted, surrounded
by other under-predicted municipalities), 554 form "Low-Low"
clusters, and 6{,}418 show no significant local pattern.

An explicit, honest limitation belongs here: explicit spatial
autoregressive models (SAR/CAR) were attempted as the natural next
step, but did not converge within a practical runtime on this
dataset size and were abandoned rather than reported with
unreliable numbers. This is stated openly, not hidden.

\subsection*{Step 5 --- Is the model statistically trustworthy, not
just a lucky R$^2$?}
A battery of additional diagnostics, summarised briefly because they
belong to the report rather than the poster: calibration by decile
(Pearson/Spearman/Lin's CCC $= 0.885/0.890/0.879$), an empirical 95\%
prediction interval of width 7.79 IFC points, AIC/BIC model
comparison confirming the nested model is preferred, and residual
normality diagnostics (skewness $0.35$, excess kurtosis $2.31$,
Shapiro--Wilk $W = 0.982$) showing mild but real deviation from
normality, fully documented with QQ plots in the dedicated
parsimonious-extension report.

\subsection*{Step 6 --- Is there a fundamentally different way to
reach the same accuracy?}
Responding directly to the instructor's remark ``PCR, not PCA'':
\begin{itemize}[leftmargin=2em]
  \item \textbf{PCR on the 326 GRINS features}: best R$^2 = 0.787$
  at $K=100$ components --- below both the parsimonious model and
  the benchmark. Principal components are chosen to maximise
  variance in $X$, not predictive power for $Y$; the directions do
  not coincide here.
  \item \textbf{PCR on the 12 raw IFC sub-indicators} (the same
  components used by the descriptive PCA cascade): R$^2 = 0.987$ at
  $K=12$, i.e. using \emph{all} the information. The 1.3\% residual
  gap is informative, not noise: it reveals that the official IFC
  construction includes a non-linear step (almost certainly a
  percentile-rank transform) beyond a linear weighted aggregation of
  its 12 components.
  \item \textbf{PLS vs PCR, head-to-head}, same 326-feature matrix:
  PLS (which chooses components by covariance with $Y$, not variance
  of $X$ alone) dominates PCR by about 5 percentage points of
  R$^2$ at every component count. PLS at $K=75$ reaches
  R$^2=0.825$, matching the 191-covariate LASSO benchmark
  ($0.826$) almost exactly.
\end{itemize}
Conclusion of this step: dimensionality reduction \emph{can} match
LASSO's accuracy, but only by trading away interpretability ---
75 latent components have no individual substantive meaning, while
LASSO's 191 (or the parsimonious 25) are named, signed, readable
variables. LASSO is retained as the project's primary deliverable on
interpretability grounds, not because it is more accurate.

\section{The five substantive drivers, explained}

Standardised effect sizes ($\beta \cdot \mathrm{SD}$): the change in
IFC per one standard deviation increase in the predictor, holding
the rest of the model fixed.

\begin{center}
\begin{tabular}{@{}cllr@{}}
\toprule
Rank & Variable & Direction & $\beta \cdot \mathrm{SD}$ \\
\midrule
1 & Middle-income taxpayer mass        & less fragile & $-1.03$ \\
2 & Total taxpayer density             & less fragile & $-0.78$ \\
3 & Low-income declared income mass    & more fragile  & $+0.57$ \\
4 & Pension share of income            & more fragile  & $+0.56$ \\
5 & Small construction firms density   & less fragile  & $-0.53$ \\
\bottomrule
\end{tabular}
\end{center}

Reading: fragility is a balance between \emph{economic dynamism}
(a broad, present middle class; a dense taxpayer base; a fabric of
small construction enterprises) and \emph{economic distress}
(concentration of low-income declared earners; heavy dependence on
pension income rather than wages). None of the top five drivers is
a demographic or environmental variable --- the strongest signal is
unambiguously fiscal/economic.

\section{What was deliberately left out of the poster, and why}

\begin{itemize}[leftmargin=2em]
  \item \textbf{Descriptive PCA on GRINS} and \textbf{k-means
  clustering on the 12 IFC sub-components} (a teammate's
  contribution): these are exploratory/contextual analyses that
  confirm fragility is geographically structured \emph{before} any
  model is fit (e.g.\ a cross-tab against the independent GRINS
  urbanisation taxonomy shows Inner Italy concentrated in deciles
  9--10, Metropolitan Italy in deciles 1--2). They motivate the
  geographic random effects but are not part of the predictive
  pipeline, so they live in the introduction/context, not in the
  main results.
  \item \textbf{Residual normality QQ plots and Box--Cox profiling}:
  real diagnostics, fully reported in the parsimonious-extension
  report, but too technical to read at poster distance in the few
  seconds a passer-by spends on any one panel.
  \item \textbf{The full 191-covariate coefficient table and the
  complete VIF-iteration log}: available in the appendix of the
  benchmark report; listing them on the poster would not add
  information a viewer could absorb on the spot.
\end{itemize}

\section{The take-home conclusion, in three points}

\begin{enumerate}[leftmargin=2em]
  \item \textbf{GRINS predicts IFC well}: R$^2 \approx 0.82$
  out-of-sample with only 25 interpretable covariates plus a
  geographic hierarchy --- within one percentage point of a
  191-covariate benchmark and of a 75-component PLS model.
  \item \textbf{The story is socio-economic, not demographic or
  environmental}: middle-class mass, taxpayer density and
  small-business fabric reduce fragility; low-income concentration
  and pension dependency increase it.
  \item \textbf{Geography matters beyond what the indicators
  capture}: a 36\% intra-class correlation plus a persistent,
  locally-confirmed Moran's $I = 0.234$ after controlling for both
  covariates and administrative geography show that explicit
  spatial models (SAR/CAR) are the natural, currently unfinished,
  next step.
\end{enumerate}

\section{Anticipated questions and prepared answers}

\textbf{Q: Why LASSO and not Ridge or plain stepwise selection?}\\
LASSO performs variable selection (sets coefficients exactly to
zero), which Ridge does not; this is essential because the goal is
an interpretable, sparse model, not just a low-error one. An Elastic
Net cross-check across $\alpha \in \{0.1, 0.25, 0.5, 0.75, 1.0\}$
showed CV error varies by less than 1\% across the whole range,
confirming the choice of $\alpha=1$ (pure LASSO) is not an arbitrary
or fragile decision.

\textbf{Q: What is the difference between a ``feature'' and a
``covariate'' in your slides?}\\
A feature is a candidate variable created during engineering (326 of
them), before any selection. A covariate is a feature that has
actually entered a fitted regression model (191 after LASSO, 25
after parsimony refinement). The words track two different stages
of the same pipeline, not two different objects.

\textbf{Q: Why cross-validate by municipality rather than by row?}\\
Each municipality appears twice (2019 and 2021). Splitting by row
would let one year of a municipality train the model while the
other year of the \emph{same} municipality tests it, leaking
information and inflating R$^2$ optimistically. Splitting whole
municipalities into folds removes this leakage entirely.

\textbf{Q: Why exclude total surface area?}\\
Per direct instructor feedback: once every other covariate is
expressed as a per-capita density or rate, surface area is a static
scale descriptor, not a fragility driver in this specification ---
keeping it would have reintroduced exactly the kind of
non-substantive, mechanical covariate the parsimonious model is
designed to avoid.

\textbf{Q: Why does PCR on the 12 official IFC sub-indicators reach
only R$^2=0.987$, not 1.0, if those are literally the index's own
components?}\\
Because PCR (and PCA) can only recover \emph{linear} combinations of
the input variables. The residual 1.3\% indicates the official IFC
construction includes a non-linear step --- most plausibly a
percentile-rank transform applied at some stage --- that a linear
projection cannot fully reproduce. This was confirmed after fixing a
genuine data-quality issue (duplicate municipality--year rows in the
official spreadsheet, which had been silently inflating the
apparent error to 1.8\% before deduplication).

\textbf{Q: Why does PLS beat PCR so consistently?}\\
PCR chooses components to maximise variance within $X$ alone; PLS
chooses components to maximise covariance between $X$ and $Y$. When
the directions of maximum variance and maximum predictive relevance
do not coincide --- which the jump observed at very low component
counts in the PCR trade-off curve suggests is the case here --- PLS
has a structural advantage that grows with the dimensionality
mismatch.

\textbf{Q: If PLS matches LASSO's accuracy, why is LASSO still your
main deliverable?}\\
Accuracy is not the only criterion the project cares about. A PLS
component is a linear combination of arbitrarily many original
variables with no individually interpretable meaning; a LASSO
coefficient is attached to one named, signed variable a policymaker
can act on. The project's brief explicitly asks for an
\emph{interpretable} benchmark, so LASSO is preferred at equal
accuracy.

\textbf{Q: Why was the explicit spatial model (SAR/CAR) not
completed?}\\
It was attempted in earnest; the model failed to converge within a
practical runtime on a dataset of this size (15{,}782 panel
observations with a dense municipal-level spatial weight matrix).
Rather than report unreliable coefficients from a non-converged fit,
the attempt is documented as an honest limitation and a clearly
identified next step, supported by the Moran's $I$ / LISA evidence
that such a model is in fact warranted.

\textbf{Q: What does the 36\% intra-class correlation actually
mean?}\\
After the 25 fixed-effect covariates have done their work, the
\emph{remaining} variance in IFC splits roughly 36\% to ``which
region/province a municipality is in'' (a label, not a measured
covariate) and 64\% to everything else (idiosyncratic municipal
variation, measurement noise, and the spatially-correlated component
later confirmed by Moran's $I$/LISA).

\textbf{Q: How confident should we be in the top-5 drivers — could
they be an artefact of multicollinearity?}\\
The parsimonious selection enforces VIF $<10$ for every retained
covariate (maximum observed VIF is $6.4$), specifically to rule out
this concern; all five top drivers remain significant at
$p<10^{-6}$ even after the iterative replacement step removed
collinear competitors.

\end{document}
